\documentclass[12pt,a4paper]{article} %12 кегль, формат бумаги А4, тип статья

\usepackage[T2A]{fontenc} %поддержка кириллицы
\usepackage[utf8]{inputenc} %кодировка символов юникод
\usepackage[english,russian]{babel} %список рабочих языков
\usepackage{cmap} %для поиска в pdf
\usepackage[left=2cm,right=2cm,top=2cm,bottom=2cm]{geometry} %размеры отступов от краев страницы
\usepackage{graphicx} %возможность работать с графикой
\usepackage{enumitem} %для работы со списками
\usepackage{float} %для точного размещения таблиц и графики
\usepackage{caption} %для создания подписей таблиц и графики
\usepackage{amsmath,amssymb,amsthm,amsfonts} %штуки для математики
\usepackage{extarrows} %чтобы подписи над и под нестандартными стрелками писать короче в коде
\usepackage{makeidx} %собирает ключевые слова для предметного указателя
\usepackage[hidelinks]{hyperref} %убирает выделение ссылок; добавляет гиперссылки в оглавление, перекрестные ссылки и URL
\usepackage{setspace} %устанавливает отступы
\linespread{1.5} %длина межстрочного интервала
\vspace{125mm} %длина абзацного отступа
\usepackage{titlesec}
\titleformat{\section}{\normalfont\Large\bfseries\centering}{}{0pt}{} % эти две строчки убирают циферки перед названиями разделов в основном тексте и выравнивают названия секций по центру
\usepackage{paratype} %чтобы точно был поиск и шрифт хороший

\author{Автор конспекта: Краснов Борис} %автор
\title{Жеданов А.С.: Введение в специальные функции} %название
\date{Осень 2026} %дата создания документа

\newtheorem{theorem}{Теорема}
\newtheorem{lemma}{Лемма}
\newtheorem{definition}{Определение}
\newtheorem{proposition}{Предложение}
\newtheorem{consequence}{Следствие}
\newcommand{\R}{\mathbb{R}}
\newcommand{\Co}{\mathbb{C}}
\newcommand{\E}{\mathbb{E}}
\newcommand{\D}{\mathbb{D}}
\newcommand{\PR}{\mathbb{P}}
\newcommand{\charf}[2][t]{\varphi_{#2}(#1)}
\newcommand{\F}{\mathcal{F}}
\newcommand{\M}{\mathfrak{M}}

\begin{document}
	\maketitle
	\tableofcontents
	\clearpage
	
	\section{Введение}
	Некоторые методы введения специальных функций:
	\begin{enumerate}
		\item Можно рассматривать интегралы с переменным верхним пределом от чего-то полезного, которые не выражаются в элементарных функциях. Например тот же интгральный синус или логарифм и прочее.
		\item При решении дифференциальных уравнений тоже возникают специальные функции, например при разделении переменных:
		$$\triangle F(x,y,z)=\lambda F(x,y,z)$$
		Мы хотим представить $F=f_1(u)f_2(v)f_3(w)$, то есть надо новые переменные выражать через старые.
	\end{enumerate}
	
	Из истории: люди рассматривали функции вида:
	$$\int_0^x\frac{dt}{\sqrt{\alpha t^4+\beta t^2+1}}$$
	Подобного рода функции классифицировал Лежандр. Позже Якоби и Абель задались вопросом о том, какие функции являются обратными к этим. Так получились эллиптические функции.
	
	Например известный нам эллиптический синус $Sn(z)$ - двоякопереодичная функции на $\Co$ (то есть живет по факту на торе) является обратной к
	$$\int_0^z\frac{dt}{\sqrt{(1-t^2)(1-\kappa^2t^2)}},\quad\kappa\in(0,1)$$
	
	Для них например выполнена теорема сложения, а именно:
	$$\forall v,w\,\exists\,\Phi(z_1,z_2,z_3)-polynom\,\forall u\,\,\Phi(f(u),f(u+v),f(u+w))=0$$
	
	Также было что-то сказано про ортогональные полиномы (я не понял к чему, разве что они выражаются через гипергеометрическую функцию Гаусса о которой написано ниже).
	$$\int_a^bP_n(x)P_m(x)w(x)dx=\delta_{n,m},$$
	где $w(x)$ некоторая фиксированная весовая функция, $P_k(x)$ полиномы степени не выше $k$.
	
	Известные примеры: если рассмотреть отрезок $[-1,1]$ с единичной весовой функцией, то получим полномы Лежандра, а если рассмотреть весовую функцию
	$$w(x)=(1-x)^a(1+x)^b$$
	то получим полиномы Якоби.
	
	Гипергеометрическая функция Гаусса:
	$$_2F_1(a,b;c;z):=\sum_{k=0}^{+\infty}A_nz^n,\quad A_n:=\frac{(a)_n(b)_n}{n!(c)_n},\quad (a)_n:=a(a+1)(a+2)\cdots(a+n-1)$$
	Эта функция является решением ДУ:
	$$z(1-z)F^{''}(z)+(\alpha z+\beta)F^{'}(z)=\gamma F(z)$$
	Также она хороша тем, что через нее выражается много элементарных.
	
	Ну и пару слов о том, почему не любая функция из не элементарных не является специальной (да, тут бы еще знать строгое определение этих классов, но это же введение, так что замнем для ясности). Например можно посмотреть на функцию Дирихле и через пару секунд забыть, так как нет у нее никаких свойств. С другой стороны есть гамма функция, которая очень естественно задается, имеет кучу рекурентных формул, имеет понятную ассимптотику и много где возникает.
\end{document}